\section{Metodología de Proyección en las Gráficas Experimentales}
\label{ap:proyeccion_graficas}

Las gráficas de la evaluación experimental distinguen, para cada caso analizado, dos
magnitudes: el tiempo total de ejecución y la memoria máxima consumida. En todas ellas, el
eje \(x\) representa el número de qubits \(n\), mientras que el eje \(y\) se presenta en escala
logarítmica. Esta elección responde a que tanto el tamaño del vector de estado como el costo
de simularlo tienden a crecer de manera multiplicativa a medida que aumenta \(n\). Usar una
escala logarítmica permite, por tanto, comparar en una misma figura valores que difieren por
varios órdenes de magnitud sin perder legibilidad.

En cada curva, las líneas continuas corresponden a mediciones obtenidas directamente en los
experimentos realizados. A partir del último punto medido, las líneas discontinuas muestran
una proyección del comportamiento esperado para tamaños no evaluados de forma directa. La
banda sombreada indica el intervalo de predicción del \(95\%\), y la línea vertical punteada
marca precisamente el límite entre la región observada y la región proyectada. De este modo,
la figura permite distinguir con claridad qué parte de la curva proviene de datos medidos y
cuál debe interpretarse como una estimación.

\subsection*{Modelo de Proyección}

La proyección se construyó de manera independiente para cada combinación de experimento,
estrategia de almacenamiento y métrica. Si \(y_i\) denota el tiempo o la memoria medidos
para un número de qubits \(n_i\), el primer paso consiste en transformar la variable de
respuesta al espacio logarítmico:
\[
z_i = \log(y_i).
\]

Esta transformación se utiliza porque el crecimiento de tiempo y memoria en simulación
cuántica de estado completo puede aproximarse razonablemente mediante una tendencia
exponencial respecto al número de qubits. Al pasar al espacio logarítmico, esa relación se
vuelve aproximadamente lineal, lo que permite ajustarla con un modelo más simple e
interpretable. Además, al regresar después a la escala original mediante la función
\(\exp(\cdot)\), las predicciones conservan valores positivos y mantienen una interpretación
coherente con magnitudes como tiempo y memoria.

En los casos en que una misma figura resume varias variantes de un experimento, como ocurre
con Grover para los objetivos ubicados en \(N/8\), \(N/2\) y \(7N/8\), primero se obtiene un
valor representativo por qubit y estrategia mediante la media geométrica. Esta elección es
consistente con el ajuste en espacio logarítmico, ya que equivale a promediar los valores
transformados y luego volver a la escala original.

Sobre los valores así obtenidos se ajusta una regresión lineal de la forma
\[
z_i = \beta_0 + \beta_1 n_i + \varepsilon_i,
\]
donde \(\varepsilon_i\) recoge la variación que no queda explicada por la tendencia
aproximadamente exponencial del fenómeno.

La estimación de los parámetros se realiza por mínimos cuadrados ordinarios. Si \(X\) es la
matriz de diseño, formada por una columna de unos y una columna con los valores de \(n_i\),
entonces
\[
\hat{\beta} = (X^\top X)^{-1} X^\top z.
\]

Para un nuevo tamaño \(n_0\), con \(x_0 = [1, n_0]^\top\), la predicción en espacio
logarítmico viene dada por
\[
\hat{z}_0 = x_0^\top \hat{\beta}.
\]
Finalmente, el valor proyectado que se muestra en la gráfica se obtiene al volver a la escala
original:
\[
\hat{y}_0 = \exp(\hat{z}_0).
\]

\subsection*{Estimación de la Incertidumbre}

La incertidumbre asociada a la proyección también se calcula en espacio logarítmico. Una vez
ajustado el modelo, la varianza residual se estima como
\[
s^2 = \frac{1}{m-p}\sum_{i=1}^{m}(z_i-\hat{z}_i)^2,
\]
donde \(m\) es el número de observaciones utilizadas en el ajuste y \(p=2\) corresponde al
número de parámetros del modelo. En las curvas presentadas en este trabajo se emplean seis
mediciones por estrategia, correspondientes a \(q20\) hasta \(q25\), por lo que los grados de
libertad son \(m-p=4\).

Para un nuevo valor \(n_0\), el error estándar de predicción se calcula como
\[
s_{\mathrm{pred}}(n_0)
=
s\sqrt{1+x_0^\top (X^\top X)^{-1}x_0}.
\]

Esta expresión incorpora dos fuentes de incertidumbre. Por un lado, el término
\[
x_0^\top (X^\top X)^{-1}x_0
\]
refleja que la proyección se vuelve menos precisa a medida que el punto estimado se aleja de la región en la que se cuenta con mediciones. Por otro, el
término adicional \(1\) indica que no se está construyendo solo un intervalo de confianza para
la media de la tendencia, sino un intervalo de predicción para una posible observación futura.

Con ello, el intervalo de predicción del \(95\%\) en espacio logarítmico se expresa como
\[
\hat{z}_0
\pm
t_{0.975,\,m-p}\,s_{\mathrm{pred}}(n_0),
\]
donde \(t_{0.975,\,m-p}\) es el valor crítico de la distribución \(t\) de Student con
\(m-p\) grados de libertad.

Para representar este intervalo en la escala original, se aplica la función exponencial a sus
dos extremos:
\[
\left[
\exp\!\left(\hat{z}_0-t_{0.975,\,m-p}s_{\mathrm{pred}}(n_0)\right),
\exp\!\left(\hat{z}_0+t_{0.975,\,m-p}s_{\mathrm{pred}}(n_0)\right)
\right].
\]

Se emplea la distribución \(t\) de Student, y no una aproximación normal, porque cada curva
se ajusta con un número reducido de observaciones. Esta elección permite incorporar de forma
más adecuada la incertidumbre asociada a la estimación de la varianza residual. Del mismo
modo, se utiliza un intervalo de predicción y no únicamente un intervalo de confianza, ya que
la banda sombreada busca indicar un rango plausible para futuras ejecuciones y no solo la
incertidumbre de la tendencia media ajustada.
